% CP-VI-0001
% target_chapter: VI/08 — Basic Open Sets
% source_unit_id: IMP-DL-0E3E7EC766-P001
% source: Downloads/theory-of-algebraic-geometry-1.tex L203-L357
% solution_provenance: SOURCE_EXPLICIT_SOLUTION

\subsection*{Problem 1 - Solutions: Basic Open Sets, Spectra of Polynomial Rings, and Disconnected Spectra}
	
	Let \(A\) be a ring. Show that the sets
	\[
	D(f)=\{\mathfrak p\in \operatorname{Spec} A \mid f\notin \mathfrak p\},
	\]
	with \(f\) ranging over all elements of \(A\), form a basis of the topology on
	\(\operatorname{Spec} A\). Show that if \(D(f)=\varnothing\), then \(f\) is nilpotent, i.e. there
	exists \(n>0\) such that \(f^n=0\).
	
	\bigskip
	
	\textbf{Solution.}
	
	Recall that the Zariski closed subsets of \(\operatorname{Spec} A\) are sets of the form
	\[
	V(I)=\{\mathfrak p\in \operatorname{Spec} A \mid I\subseteq \mathfrak p\},
	\]
	where \(I\subseteq A\) is an ideal.
	
	For a single element \(f\in A\), we have
	\[
	V((f))=\{\mathfrak p\in \operatorname{Spec} A \mid f\in \mathfrak p\}.
	\]
	Therefore
	\[
	D(f)=\{\mathfrak p\in \operatorname{Spec} A \mid f\notin \mathfrak p\}
	=\operatorname{Spec} A\setminus V((f)).
	\]
	Hence every \(D(f)\) is open.
	
	We now show that the sets \(D(f)\) form a basis.
	
	First,
	\[
	D(1)=\operatorname{Spec} A,
	\]
	because no prime ideal contains \(1\). Thus the sets \(D(f)\) cover \(\operatorname{Spec} A\).
	
	Second, for any \(f,g\in A\), we have
	\[
	D(f)\cap D(g)=D(fg).
	\]
	Indeed,
	\[
	\mathfrak p\in D(f)\cap D(g)
	\]
	if and only if
	\[
	f\notin \mathfrak p
	\quad\text{and}\quad
	g\notin \mathfrak p.
	\]
	Since \(\mathfrak p\) is prime, this is equivalent to
	\[
	fg\notin \mathfrak p.
	\]
	Therefore
	\[
	D(f)\cap D(g)=D(fg).
	\]
	
	Now let \(U\subseteq \operatorname{Spec} A\) be open. Then
	\[
	U=\operatorname{Spec} A\setminus V(I)
	\]
	for some ideal \(I\subseteq A\). We claim that
	\[
	U=\bigcup_{f\in I}D(f).
	\]
	Indeed,
	\[
	\mathfrak p\in \operatorname{Spec} A\setminus V(I)
	\]
	if and only if
	\[
	I\not\subseteq \mathfrak p.
	\]
	This means that there exists \(f\in I\) such that
	\[
	f\notin \mathfrak p.
	\]
	Equivalently,
	\[
	\mathfrak p\in D(f)
	\]
	for some \(f\in I\). Hence
	\[
	\operatorname{Spec} A\setminus V(I)=\bigcup_{f\in I}D(f).
	\]
	Therefore every open set is a union of basic opens \(D(f)\), so the sets \(D(f)\)
	form a basis for the topology on \(\operatorname{Spec} A\).
	
	\bigskip
	
	Now we prove the nilpotence statement.
	
	We first recall the classical identity
	\[
	\sqrt{(0)}
	=
	\bigcap_{\mathfrak p\in \operatorname{Spec} A}\mathfrak p.
	\]
	The ideal \(\sqrt{(0)}\) is called the nilradical of \(A\), and it consists
	exactly of the nilpotent elements of \(A\).
	
	Suppose first that
	\[
	D(f)=\varnothing.
	\]
	Then there is no prime ideal \(\mathfrak p\) such that \(f\notin \mathfrak p\).
	Equivalently,
	\[
	f\in \mathfrak p
	\]
	for every prime ideal \(\mathfrak p\subseteq A\). Hence
	\[
	f\in \bigcap_{\mathfrak p\in \operatorname{Spec} A}\mathfrak p.
	\]
	Using the identity above, we get
	\[
	f\in \sqrt{(0)}.
	\]
	Therefore \(f\) is nilpotent. Hence there exists \(n>0\) such that
	\[
	f^n=0.
	\]
	
	Conversely, suppose \(f\) is nilpotent. Then there exists \(n>0\) such that
	\[
	f^n=0.
	\]
	Let \(\mathfrak p\in \operatorname{Spec} A\). Since \(0\in \mathfrak p\), we have
	\[
	f^n\in \mathfrak p.
	\]
	Because \(\mathfrak p\) is prime, this implies
	\[
	f\in \mathfrak p.
	\]
	Thus every prime ideal contains \(f\). Therefore there is no prime ideal avoiding
	\(f\), and so
	\[
	D(f)=\varnothing.
	\]
	
	Thus
	\[
	\boxed{
		D(f)=\varnothing
		\iff
		f\text{ is nilpotent}.
	}
	\]
